% Gaya bersama untuk kedua lembar kuesioner penelitian Bab 16.
\usepackage[a4paper, top=1.6cm, bottom=1.6cm, left=2cm, right=2cm]{geometry}
\usepackage{fontspec}
\setmainfont[
  Path = ../fonts/,
  UprightFont = TitilliumWeb-Regular.ttf,
  BoldFont = TitilliumWeb-Bold.ttf,
  ItalicFont = TitilliumWeb-Italic.ttf,
  BoldItalicFont = TitilliumWeb-BoldItalic.ttf
]{TitilliumWeb}
\usepackage[bahasai]{babel}
\usepackage{xcolor}
\usepackage{graphicx}
\usepackage{tabularx}
\usepackage{array}
\usepackage{tikz}

\definecolor{praditagreen}{RGB}{37, 113, 71}
\setlength{\parindent}{0pt}
\setlength{\parskip}{0.5em}
\pagestyle{empty}

% Bagian yang harus diisi peneliti sebelum lembar dicetak.
\newcommand{\isian}[1]{\textcolor{red}{[#1]}}

% Kepala lembar: logo, judul penelitian, dan nomor lembar.
\newcommand{\kepala}[1]{%
  \begin{minipage}[c]{0.22\textwidth}
    \includegraphics[width=\linewidth]{../../figures/images/logo_universitas.png}
  \end{minipage}\hfill
  \begin{minipage}[c]{0.74\textwidth}
    \raggedleft
    {\Large\bfseries\color{praditagreen} Kuesioner Penelitian}\\[2pt]
    {\large Belajar Pemrograman dengan Bantuan AI}\\[2pt]
    {\bfseries #1}
  \end{minipage}
  \par\vspace{2pt}{\color{praditagreen}\rule{\textwidth}{1pt}}\par
}

\newcommand{\bagian}[1]{\par\medskip{\large\bfseries\color{praditagreen}#1}\par}

% Kotak dan lingkaran digambar dengan TikZ agar tidak bergantung pada font simbol.
\newcommand{\kotakpilih}{\tikz[baseline=-0.15em]\draw (0,0) rectangle (0.3,0.3);}
\newcommand{\lingkaran}{\tikz[baseline=-0.6ex]\draw (0,0) circle (0.17);}

% Pilihan jawaban: kotak tidak terpisah dari labelnya saat baris dipatah.
\newcommand{\opsi}[1]{\mbox{\kotakpilih~#1}\hspace{1em plus 0.5em}}
\newenvironment{barispilihan}{\par\raggedright}{\par}

% NIM ditulis di kedua lembar untuk memasangkannya. NIM diganti kode acak
% sebelum analisis dan tidak dimuat dalam publikasi.
\newcommand{\kotaknim}{%
  Tuliskan NIM Anda, satu angka di setiap kotak.

  \begin{tabular}{|*{10}{p{0.45cm}|}}
    \hline
    \rule{0pt}{0.8cm} & & & & & & & & & \\
    \hline
  \end{tabular}
}

% Tabel butir dengan lima pilihan lingkaran. Setiap baris diakhiri \pilihan.
\newenvironment{tabelskala}{%
  \renewcommand{\arraystretch}{1.45}%
  \tabularx{\textwidth}{>{\raggedright\arraybackslash}p{0.8cm}X*{5}{>{\centering\arraybackslash}p{0.65cm}}}
  \hline
  \textbf{No} & \textbf{Pernyataan} & \textbf{1} & \textbf{2} & \textbf{3} & \textbf{4} & \textbf{5} \\
  \hline
}{\endtabularx}
\newcommand{\pilihan}{& \lingkaran & \lingkaran & \lingkaran & \lingkaran & \lingkaran \\ \hline}

\newcommand{\garis}{\par\vspace{0.55cm}\noindent\rule{\textwidth}{0.4pt}}
